% Options for packages loaded elsewhere
% Options for packages loaded elsewhere
\PassOptionsToPackage{unicode}{hyperref}
\PassOptionsToPackage{hyphens}{url}
%
\documentclass[
  letterpaper,
]{scrbook}
\usepackage{xcolor}
\usepackage{amsmath,amssymb}
\setcounter{secnumdepth}{-\maxdimen} % remove section numbering
\usepackage{iftex}
\ifPDFTeX
  \usepackage[T1]{fontenc}
  \usepackage[utf8]{inputenc}
  \usepackage{textcomp} % provide euro and other symbols
\else % if luatex or xetex
  \usepackage{unicode-math} % this also loads fontspec
  \defaultfontfeatures{Scale=MatchLowercase}
  \defaultfontfeatures[\rmfamily]{Ligatures=TeX,Scale=1}
\fi
\usepackage{lmodern}
\ifPDFTeX\else
  % xetex/luatex font selection
\fi
% Use upquote if available, for straight quotes in verbatim environments
\IfFileExists{upquote.sty}{\usepackage{upquote}}{}
\IfFileExists{microtype.sty}{% use microtype if available
  \usepackage[]{microtype}
  \UseMicrotypeSet[protrusion]{basicmath} % disable protrusion for tt fonts
}{}
\makeatletter
\@ifundefined{KOMAClassName}{% if non-KOMA class
  \IfFileExists{parskip.sty}{%
    \usepackage{parskip}
  }{% else
    \setlength{\parindent}{0pt}
    \setlength{\parskip}{6pt plus 2pt minus 1pt}}
}{% if KOMA class
  \KOMAoptions{parskip=half}}
\makeatother
% Make \paragraph and \subparagraph free-standing
\makeatletter
\ifx\paragraph\undefined\else
  \let\oldparagraph\paragraph
  \renewcommand{\paragraph}{
    \@ifstar
      \xxxParagraphStar
      \xxxParagraphNoStar
  }
  \newcommand{\xxxParagraphStar}[1]{\oldparagraph*{#1}\mbox{}}
  \newcommand{\xxxParagraphNoStar}[1]{\oldparagraph{#1}\mbox{}}
\fi
\ifx\subparagraph\undefined\else
  \let\oldsubparagraph\subparagraph
  \renewcommand{\subparagraph}{
    \@ifstar
      \xxxSubParagraphStar
      \xxxSubParagraphNoStar
  }
  \newcommand{\xxxSubParagraphStar}[1]{\oldsubparagraph*{#1}\mbox{}}
  \newcommand{\xxxSubParagraphNoStar}[1]{\oldsubparagraph{#1}\mbox{}}
\fi
\makeatother

\usepackage{color}
\usepackage{fancyvrb}
\newcommand{\VerbBar}{|}
\newcommand{\VERB}{\Verb[commandchars=\\\{\}]}
\DefineVerbatimEnvironment{Highlighting}{Verbatim}{commandchars=\\\{\}}
% Add ',fontsize=\small' for more characters per line
\usepackage{framed}
\definecolor{shadecolor}{RGB}{241,243,245}
\newenvironment{Shaded}{\begin{snugshade}}{\end{snugshade}}
\newcommand{\AlertTok}[1]{\textcolor[rgb]{0.68,0.00,0.00}{#1}}
\newcommand{\AnnotationTok}[1]{\textcolor[rgb]{0.37,0.37,0.37}{#1}}
\newcommand{\AttributeTok}[1]{\textcolor[rgb]{0.40,0.45,0.13}{#1}}
\newcommand{\BaseNTok}[1]{\textcolor[rgb]{0.68,0.00,0.00}{#1}}
\newcommand{\BuiltInTok}[1]{\textcolor[rgb]{0.00,0.23,0.31}{#1}}
\newcommand{\CharTok}[1]{\textcolor[rgb]{0.13,0.47,0.30}{#1}}
\newcommand{\CommentTok}[1]{\textcolor[rgb]{0.37,0.37,0.37}{#1}}
\newcommand{\CommentVarTok}[1]{\textcolor[rgb]{0.37,0.37,0.37}{\textit{#1}}}
\newcommand{\ConstantTok}[1]{\textcolor[rgb]{0.56,0.35,0.01}{#1}}
\newcommand{\ControlFlowTok}[1]{\textcolor[rgb]{0.00,0.23,0.31}{\textbf{#1}}}
\newcommand{\DataTypeTok}[1]{\textcolor[rgb]{0.68,0.00,0.00}{#1}}
\newcommand{\DecValTok}[1]{\textcolor[rgb]{0.68,0.00,0.00}{#1}}
\newcommand{\DocumentationTok}[1]{\textcolor[rgb]{0.37,0.37,0.37}{\textit{#1}}}
\newcommand{\ErrorTok}[1]{\textcolor[rgb]{0.68,0.00,0.00}{#1}}
\newcommand{\ExtensionTok}[1]{\textcolor[rgb]{0.00,0.23,0.31}{#1}}
\newcommand{\FloatTok}[1]{\textcolor[rgb]{0.68,0.00,0.00}{#1}}
\newcommand{\FunctionTok}[1]{\textcolor[rgb]{0.28,0.35,0.67}{#1}}
\newcommand{\ImportTok}[1]{\textcolor[rgb]{0.00,0.46,0.62}{#1}}
\newcommand{\InformationTok}[1]{\textcolor[rgb]{0.37,0.37,0.37}{#1}}
\newcommand{\KeywordTok}[1]{\textcolor[rgb]{0.00,0.23,0.31}{\textbf{#1}}}
\newcommand{\NormalTok}[1]{\textcolor[rgb]{0.00,0.23,0.31}{#1}}
\newcommand{\OperatorTok}[1]{\textcolor[rgb]{0.37,0.37,0.37}{#1}}
\newcommand{\OtherTok}[1]{\textcolor[rgb]{0.00,0.23,0.31}{#1}}
\newcommand{\PreprocessorTok}[1]{\textcolor[rgb]{0.68,0.00,0.00}{#1}}
\newcommand{\RegionMarkerTok}[1]{\textcolor[rgb]{0.00,0.23,0.31}{#1}}
\newcommand{\SpecialCharTok}[1]{\textcolor[rgb]{0.37,0.37,0.37}{#1}}
\newcommand{\SpecialStringTok}[1]{\textcolor[rgb]{0.13,0.47,0.30}{#1}}
\newcommand{\StringTok}[1]{\textcolor[rgb]{0.13,0.47,0.30}{#1}}
\newcommand{\VariableTok}[1]{\textcolor[rgb]{0.07,0.07,0.07}{#1}}
\newcommand{\VerbatimStringTok}[1]{\textcolor[rgb]{0.13,0.47,0.30}{#1}}
\newcommand{\WarningTok}[1]{\textcolor[rgb]{0.37,0.37,0.37}{\textit{#1}}}

\usepackage{longtable,booktabs,array}
\newcounter{none} % for unnumbered tables
\usepackage{calc} % for calculating minipage widths
% Correct order of tables after \paragraph or \subparagraph
\usepackage{etoolbox}
\makeatletter
\patchcmd\longtable{\par}{\if@noskipsec\mbox{}\fi\par}{}{}
\makeatother
% Allow footnotes in longtable head/foot
\IfFileExists{footnotehyper.sty}{\usepackage{footnotehyper}}{\usepackage{footnote}}
\makesavenoteenv{longtable}
\usepackage{graphicx}
\makeatletter
\newsavebox\pandoc@box
\newcommand*\pandocbounded[1]{% scales image to fit in text height/width
  \sbox\pandoc@box{#1}%
  \Gscale@div\@tempa{\textheight}{\dimexpr\ht\pandoc@box+\dp\pandoc@box\relax}%
  \Gscale@div\@tempb{\linewidth}{\wd\pandoc@box}%
  \ifdim\@tempb\p@<\@tempa\p@\let\@tempa\@tempb\fi% select the smaller of both
  \ifdim\@tempa\p@<\p@\scalebox{\@tempa}{\usebox\pandoc@box}%
  \else\usebox{\pandoc@box}%
  \fi%
}
% Set default figure placement to htbp
\def\fps@figure{htbp}
\makeatother





\setlength{\emergencystretch}{3em} % prevent overfull lines

\providecommand{\tightlist}{%
  \setlength{\itemsep}{0pt}\setlength{\parskip}{0pt}}



 


\makeatletter
\@ifpackageloaded{tcolorbox}{}{\usepackage[skins,breakable]{tcolorbox}}
\@ifpackageloaded{fontawesome5}{}{\usepackage{fontawesome5}}
\definecolor{quarto-callout-color}{HTML}{909090}
\definecolor{quarto-callout-note-color}{HTML}{0758E5}
\definecolor{quarto-callout-important-color}{HTML}{CC1914}
\definecolor{quarto-callout-warning-color}{HTML}{EB9113}
\definecolor{quarto-callout-tip-color}{HTML}{00A047}
\definecolor{quarto-callout-caution-color}{HTML}{FC5300}
\definecolor{quarto-callout-color-frame}{HTML}{acacac}
\definecolor{quarto-callout-note-color-frame}{HTML}{4582ec}
\definecolor{quarto-callout-important-color-frame}{HTML}{d9534f}
\definecolor{quarto-callout-warning-color-frame}{HTML}{f0ad4e}
\definecolor{quarto-callout-tip-color-frame}{HTML}{02b875}
\definecolor{quarto-callout-caution-color-frame}{HTML}{fd7e14}
\makeatother
\makeatletter
\@ifpackageloaded{bookmark}{}{\usepackage{bookmark}}
\makeatother
\makeatletter
\@ifpackageloaded{caption}{}{\usepackage{caption}}
\AtBeginDocument{%
\ifdefined\contentsname
  \renewcommand*\contentsname{Table of contents}
\else
  \newcommand\contentsname{Table of contents}
\fi
\ifdefined\listfigurename
  \renewcommand*\listfigurename{List of Figures}
\else
  \newcommand\listfigurename{List of Figures}
\fi
\ifdefined\listtablename
  \renewcommand*\listtablename{List of Tables}
\else
  \newcommand\listtablename{List of Tables}
\fi
\ifdefined\figurename
  \renewcommand*\figurename{Figure}
\else
  \newcommand\figurename{Figure}
\fi
\ifdefined\tablename
  \renewcommand*\tablename{Table}
\else
  \newcommand\tablename{Table}
\fi
}
\@ifpackageloaded{float}{}{\usepackage{float}}
\floatstyle{ruled}
\@ifundefined{c@chapter}{\newfloat{codelisting}{h}{lop}}{\newfloat{codelisting}{h}{lop}[chapter]}
\floatname{codelisting}{Listing}
\newcommand*\listoflistings{\listof{codelisting}{List of Listings}}
\makeatother
\makeatletter
\makeatother
\makeatletter
\@ifpackageloaded{caption}{}{\usepackage{caption}}
\@ifpackageloaded{subcaption}{}{\usepackage{subcaption}}
\makeatother
\usepackage{bookmark}
\IfFileExists{xurl.sty}{\usepackage{xurl}}{} % add URL line breaks if available
\urlstyle{same}
\hypersetup{
  pdftitle={从样地到结论：保护生物学数据实战},
  pdfauthor={管振华},
  hidelinks,
  pdfcreator={LaTeX via pandoc}}


\title{从样地到结论：保护生物学数据实战}
\usepackage{etoolbox}
\makeatletter
\providecommand{\subtitle}[1]{% add subtitle to \maketitle
  \apptocmd{\@title}{\par {\large #1 \par}}{}{}
}
\makeatother
\subtitle{基于 R 语言的生物实验设计与统计分析}
\author{管振华}
\date{2026-01-01}
\begin{document}
\frontmatter
\maketitle

\renewcommand*\contentsname{Table of contents}
{
\setcounter{tocdepth}{2}
\tableofcontents
}

\mainmatter
\bookmarksetup{startatroot}

\chapter{从样地到结论：保护生物学数据实战}\label{ux4eceux6837ux5730ux5230ux7ed3ux8bbaux4fddux62a4ux751fux7269ux5b66ux6570ux636eux5b9eux6218}

基于 R 语言的生物实验设计与统计分析

\hfill\break

\section{欢迎}\label{ux6b22ux8fce}

本书是一本面向野生动物保护、森林保护、植物保护、林学与生物多样性专业学生的案例驱动统计教材。

\subsection{这本书是什么}\label{ux8fd9ux672cux4e66ux662fux4ec0ux4e48}

\begin{itemize}
\tightlist
\item
  \textbf{保护问题先行}：每章从一个真实的保护生物学问题出发
\item
  \textbf{统计判断为核心}：重点不是背公式，而是''这个问题该用什么方法''
\item
  \textbf{R 语言实操}：所有分析代码可运行、可复现
\item
  \textbf{AI 辅助但不依赖}：教你用 AI 帮忙，同时教你检查 AI 的对错
\end{itemize}

\subsection{如何使用}\label{ux5982ux4f55ux4f7fux7528}

\begin{itemize}
\tightlist
\item
  \textbf{学生}：按章节顺序阅读，完成每章练习和 AI 核查任务
\item
  \textbf{教师}：根据课时需求选取章节，每个案例可独立教学
\item
  \textbf{保护工作者}：按专业方向或问题类型查找所需方法
\end{itemize}

\part{第一篇：从保护问题进入统计}

\chapter{第1章
统计学在保护中的角色}\label{ux7b2c1ux7ae0-ux7edfux8ba1ux5b66ux5728ux4fddux62a4ux4e2dux7684ux89d2ux8272}

\chapter{第1章
统计学在保护中的角色}\label{ux7b2c1ux7ae0-ux7edfux8ba1ux5b66ux5728ux4fddux62a4ux4e2dux7684ux89d2ux8272-1}

\section{真实保护问题}\label{ux771fux5b9eux4fddux62a4ux95eeux9898}

你管理着一片保护区，想知道：

\begin{itemize}
\tightlist
\item
  今年的病虫害是否比去年严重？
\item
  新采用的防治措施是否真的有效？
\item
  红外相机拍到的动物次数能代表真实种群数量吗？
\end{itemize}

这些问题都离不开一个工具：\textbf{统计学}。

\section{学习目标}\label{ux5b66ux4e60ux76eeux6807}

学完本章后，你应能够： 1. 区分描述、推断、预测和因果四种统计目标 2.
判断一个保护类问题是否需要统计分析 3. 理解统计学、R 和 AI 各自的责任边界
4. 建立''保护问题→数据→证据→行动''的思维框架

\section{为什么保护工作需要统计}\label{ux4e3aux4ec0ux4e48ux4fddux62a4ux5de5ux4f5cux9700ux8981ux7edfux8ba1}

保护生物学是一门''证据驱动''的学科。我们做的每一个判断------这种农药有没有用、这个物种数量是不是在下降、这片保护区要不要扩大------都需要数据支撑。

但数据不会自己说话。你需要统计来：

\begin{itemize}
\tightlist
\item
  \textbf{描述}：今年虫口密度平均多少？各林型之间差异多大？
\item
  \textbf{推断}：观察到的差异是真实存在的还是碰巧？
\item
  \textbf{预测}：如果气候变化继续，这个物种的分布区会怎么变？
\item
  \textbf{决策}：基于证据，我们应该采取什么保护行动？
\end{itemize}

\section{统计学能帮我们回答什么问题}\label{ux7edfux8ba1ux5b66ux80fdux5e2eux6211ux4eecux56deux7b54ux4ec0ux4e48ux95eeux9898}

判断以下哪些需要统计分析：

\begin{enumerate}
\def\labelenumi{\arabic{enumi}.}
\tightlist
\item
  某药剂是否降低害虫数量？ → \textbf{需要}（比较处理vs对照）
\item
  不同林型中鸟类数量是否不同？ → \textbf{需要}（多组比较）
\item
  湿度是否影响植物病害发生率？ → \textbf{需要}（相关/回归分析）
\item
  红外相机拍到的次数能否代表动物活动强度？ →
  \textbf{需要}（检测概率建模）
\end{enumerate}

结论：保护生物学中几乎所有经验性问题都需要统计分析。

\section{统计、R 与 AI
的责任边界}\label{ux7edfux8ba1r-ux4e0e-ai-ux7684ux8d23ux4efbux8fb9ux754c}

{\def\LTcaptype{none} % do not increment counter
\begin{longtable}[]{@{}lll@{}}
\toprule\noalign{}
工具 & 能做什么 & 不能做什么 \\
\midrule\noalign{}
\endhead
\bottomrule\noalign{}
\endlastfoot
统计方法 & 量化不确定性，提供证据强度 & 代替你判断''效果够不够大'' \\
R 语言 & 高效计算、可重复分析 & 替你选择正确的统计方法 \\
AI 助手 & 解释概念、推荐方法、生成代码 & 保证方法正确------你需要核查 \\
\end{longtable}
}

\section{从''是否显著''走向''下一步做什么''}\label{ux4eceux662fux5426ux663eux8457ux8d70ux5411ux4e0bux4e00ux6b65ux505aux4ec0ux4e48}

本书的核心信念：

\begin{quote}
统计不是终点，是通往保护行动的桥梁。
\end{quote}

\begin{itemize}
\tightlist
\item
  P \textless{} 0.05
  不是结论------你还需要知道效应有多大、在管理上意味着什么
\item
  P \textgreater{} 0.05
  不是失败------它告诉你当前数据不足以区分信号和噪声，可能需要更多数据或调整监测方案
\end{itemize}

\section{R 分析预览}\label{r-ux5206ux6790ux9884ux89c8}

\begin{Shaded}
\begin{Highlighting}[]
\CommentTok{\# 这是你在本书中学到的第一个 R 代码}
\CommentTok{\# 后面每章都会有可运行的示例}

\CommentTok{\# 一个简单的比较}
\NormalTok{treatment }\OtherTok{\textless{}{-}} \FunctionTok{c}\NormalTok{(}\FloatTok{3.2}\NormalTok{, }\FloatTok{2.8}\NormalTok{, }\FloatTok{3.5}\NormalTok{, }\FloatTok{2.9}\NormalTok{, }\FloatTok{3.1}\NormalTok{)  }\CommentTok{\# 处理组虫口密度}
\NormalTok{control  }\OtherTok{\textless{}{-}} \FunctionTok{c}\NormalTok{(}\FloatTok{5.8}\NormalTok{, }\FloatTok{6.1}\NormalTok{, }\FloatTok{5.5}\NormalTok{, }\FloatTok{6.3}\NormalTok{, }\FloatTok{5.7}\NormalTok{)  }\CommentTok{\# 对照组虫口密度}

\FunctionTok{mean}\NormalTok{(treatment)}
\FunctionTok{mean}\NormalTok{(control)}
\end{Highlighting}
\end{Shaded}

\section{AI 辅助任务}\label{ai-ux8f85ux52a9ux4efbux52a1}

向 AI 提问''统计学在植物保护中有什么用？``然后：

\begin{itemize}
\tightlist
\item
  检查 AI 的回答中哪些说法太笼统
\item
  找出 AI 没有提到但你学过的内容
\item
  写一句''人工核查结果''
\end{itemize}

\section{练习}\label{ux7ec3ux4e60}

\begin{enumerate}
\def\labelenumi{\arabic{enumi}.}
\tightlist
\item
  在你的专业方向中，举一个需要统计分析的例子
\item
  判断这个例子属于''描述、推断、预测、决策''中的哪一类
\end{enumerate}

\section{继续学习}\label{ux7ee7ux7eedux5b66ux4e60}

→ \href{02-variables-hypotheses.qmd}{第2章：研究问题、变量与假设}

\chapter{第2章
研究问题、变量与假设}\label{ux7b2c2ux7ae0-ux7814ux7a76ux95eeux9898ux53d8ux91cfux4e0eux5047ux8bbe}

\chapter{第2章
研究问题、变量与假设}\label{ux7b2c2ux7ae0-ux7814ux7a76ux95eeux9898ux53d8ux91cfux4e0eux5047ux8bbe-1}

\section{真实保护问题}\label{ux771fux5b9eux4fddux62a4ux95eeux9898-1}

你在评估三种药剂对森林害虫的防治效果。选了什么树？测了什么指标？怎么知道差异是药效还是碰巧？

在碰统计软件之前，你必须先回答三个问题：\textbf{研究对象是谁？变量是什么类型？你的假设是什么？}

\section{学习目标}\label{ux5b66ux4e60ux76eeux6807-1}

\begin{enumerate}
\def\labelenumi{\arabic{enumi}.}
\tightlist
\item
  将保护问题拆解为研究对象、处理因素、观测指标和研究假设
\item
  识别连续型、计数型、分类变量和比例变量
\item
  理解变量类型决定后续统计方法选择
\end{enumerate}

\section{从保护问题到统计问题}\label{ux4eceux4fddux62a4ux95eeux9898ux5230ux7edfux8ba1ux95eeux9898}

任何一个保护类研究问题都可以拆成四个要素：

{\def\LTcaptype{none} % do not increment counter
\begin{longtable}[]{@{}lll@{}}
\toprule\noalign{}
要素 & 含义 & 例子 \\
\midrule\noalign{}
\endhead
\bottomrule\noalign{}
\endlastfoot
研究对象 & 你研究的是谁 & 某森林样地中的云南松 \\
处理因素 & 你改变或比较的是什么 & 三种不同药剂 + 对照组 \\
观测指标 & 你测量的是什么 & 施药后7天的害虫死亡率 \\
研究假设 & 你预期什么结果 & 不同药剂的害虫死亡率存在差异 \\
\end{longtable}
}

\begin{tcolorbox}[enhanced jigsaw, arc=.35mm, bottomrule=.15mm, bottomtitle=1mm, breakable, colback=white, colbacktitle=quarto-callout-tip-color!10!white, colframe=quarto-callout-tip-color-frame, coltitle=black, left=2mm, leftrule=.75mm, opacityback=0, opacitybacktitle=0.6, rightrule=.15mm, title=\textcolor{quarto-callout-tip-color}{\faLightbulb}\hspace{0.5em}{课堂技巧}, titlerule=0mm, toprule=.15mm, toptitle=1mm]

拿到一段研究描述后，先用这四要素框架拆解，再考虑统计方法。很多学生跳过了这一步，直接跑t检验，结果发现变量类型根本不支持。

\end{tcolorbox}

\section{变量类型}\label{ux53d8ux91cfux7c7bux578b}

变量类型决定了你能用什么统计方法。生物统计中主要分四类：

{\def\LTcaptype{none} % do not increment counter
\begin{longtable}[]{@{}
  >{\raggedright\arraybackslash}p{(\linewidth - 6\tabcolsep) * \real{0.1875}}
  >{\raggedright\arraybackslash}p{(\linewidth - 6\tabcolsep) * \real{0.1875}}
  >{\raggedright\arraybackslash}p{(\linewidth - 6\tabcolsep) * \real{0.3438}}
  >{\raggedright\arraybackslash}p{(\linewidth - 6\tabcolsep) * \real{0.2812}}@{}}
\toprule\noalign{}
\begin{minipage}[b]{\linewidth}\raggedright
类型
\end{minipage} & \begin{minipage}[b]{\linewidth}\raggedright
定义
\end{minipage} & \begin{minipage}[b]{\linewidth}\raggedright
保护类例子
\end{minipage} & \begin{minipage}[b]{\linewidth}\raggedright
适用方法
\end{minipage} \\
\midrule\noalign{}
\endhead
\bottomrule\noalign{}
\endlastfoot
\textbf{连续型} & 可取任意数值 & 树高(m)、胸径(cm)、虫口密度 &
t检验、ANOVA、回归 \\
\textbf{计数型} & 非负整数 & 红外相机拍摄次数、叶片数 &
Poisson回归、负二项 \\
\textbf{分类变量} & 有限类别 & 栖息地类型、处理组别 & 卡方检验、ANOVA \\
\textbf{比例变量} & 0-1之间或\% & 发病率、死亡率、存活率 &
卡方检验、二项回归 \\
\end{longtable}
}

\begin{tcolorbox}[enhanced jigsaw, arc=.35mm, bottomrule=.15mm, bottomtitle=1mm, breakable, colback=white, colbacktitle=quarto-callout-warning-color!10!white, colframe=quarto-callout-warning-color-frame, coltitle=black, left=2mm, leftrule=.75mm, opacityback=0, opacitybacktitle=0.6, rightrule=.15mm, title=\textcolor{quarto-callout-warning-color}{\faExclamationTriangle}\hspace{0.5em}{关键区分}, titlerule=0mm, toprule=.15mm, toptitle=1mm]

计数型 ≠
连续型。虫口密度如果是''每株树上的虫数''，那是计数型；如果是''每株树上的虫体重量(g)``，那是连续型。

\textbf{变量类型由数据生成过程决定，不由你的测量工具决定。}

\end{tcolorbox}

\section{研究假设}\label{ux7814ux7a76ux5047ux8bbe}

研究假设是将专业问题转化为统计问题的桥梁。

\textbf{例子}：``三种药剂对害虫死亡率的影响''

\begin{itemize}
\tightlist
\item
  零假设 H₀：三种药剂的害虫死亡率没有差异
\item
  备择假设 H₁：至少有一种药剂的死亡率不同于其他
\end{itemize}

\begin{tcolorbox}[enhanced jigsaw, arc=.35mm, bottomrule=.15mm, bottomtitle=1mm, breakable, colback=white, colbacktitle=quarto-callout-important-color!10!white, colframe=quarto-callout-important-color-frame, coltitle=black, left=2mm, leftrule=.75mm, opacityback=0, opacitybacktitle=0.6, rightrule=.15mm, title=\textcolor{quarto-callout-important-color}{\faExclamation}\hspace{0.5em}{牢记}, titlerule=0mm, toprule=.15mm, toptitle=1mm]

假设必须在查看数据之前确定。先看数据再编假设，是严重的学术不端（HARKing：Hypothesizing
After Results are Known）。

\end{tcolorbox}

\section{常见错误}\label{ux5e38ux89c1ux9519ux8bef}

\begin{enumerate}
\def\labelenumi{\arabic{enumi}.}
\tightlist
\item
  \textbf{变量类型判断错误}------把计数型当连续型，用t检验代替Poisson回归
\item
  \textbf{假设不清晰就开始分析}------不知道想回答什么问题就收集数据
\item
  \textbf{混淆自变量和因变量}------``湿度影响病害''中，湿度是自变量，病害是因变量
\end{enumerate}

\section{练习}\label{ux7ec3ux4e60-1}

\begin{enumerate}
\def\labelenumi{\arabic{enumi}.}
\tightlist
\item
  判断以下变量的类型：树高、昆虫数量、栖息地类型、发病率
\item
  将''不同林型中鸟类数量是否不同''拆解为四要素
\item
  写出''湿度是否影响植物病害发生率''的零假设和备择假设
\end{enumerate}

\chapter{第3章
调查与实验设计}\label{ux7b2c3ux7ae0-ux8c03ux67e5ux4e0eux5b9eux9a8cux8bbeux8ba1}

\chapter{第3章
调查与实验设计}\label{ux7b2c3ux7ae0-ux8c03ux67e5ux4e0eux5b9eux9a8cux8bbeux8ba1-1}

\section{真实保护问题}\label{ux771fux5b9eux4fddux62a4ux95eeux9898-2}

你在评估一种新型生物农药对松材线虫病的效果。选了10棵树喷药，10棵不喷。两个月后，处理组平均病情指数2.1，对照组5.8。能说药有效吗？

\textbf{不一定。}
如果你把10棵最健康的树分到了处理组，把10棵已经病重的树分到了对照组------那差异是药效还是你分组不公平？

这就是实验设计要解决的问题。

\section{学习目标}\label{ux5b66ux4e60ux76eeux6807-2}

\begin{enumerate}
\def\labelenumi{\arabic{enumi}.}
\tightlist
\item
  理解对照、重复、随机三个实验设计基本原则
\item
  能区分生物学重复和技术重复，识别伪重复
\item
  了解完全随机设计、随机区组设计和常见保护类调查设计
\end{enumerate}

\section{三个基本原则}\label{ux4e09ux4e2aux57faux672cux539fux5219}

\subsection{1. 对照（Control）}\label{ux5bf9ux7167control}

设置未处理的参照组，排除时间、环境等无关因素的影响。

\begin{quote}
如果没有对照组，你不知道害虫是药死的还是自然死的。
\end{quote}

\subsection{2. 重复（Replication）}\label{ux91cdux590dreplication}

每个处理设置多个独立实验单位，用于估计误差。

\subsection{3. 随机（Randomization）}\label{ux968fux673arandomization}

随机分配处理到实验单位，消除系统性偏差。

\begin{quote}
不能把靠路边的样地全分给处理组，把深山里的全分给对照组。
\end{quote}

\section{生物学重复 vs
技术重复}\label{ux751fux7269ux5b66ux91cdux590d-vs-ux6280ux672fux91cdux590d}

这是生态学中最常见的统计错误来源：

{\def\LTcaptype{none} % do not increment counter
\begin{longtable}[]{@{}lll@{}}
\toprule\noalign{}
场景 & 类型 & 为什么 \\
\midrule\noalign{}
\endhead
\bottomrule\noalign{}
\endlastfoot
同一株植物测5片叶子 & 技术重复 & 5片叶子属于同一株，不是独立样本 \\
5株植物各测1片叶子 & 生物学重复 & 5株植物是独立个体 \\
同一个样方连续数5次昆虫 & 技术重复 & 同一个样方 \\
5个样方各调查一次昆虫 & 生物学重复 & 5个独立样方 \\
\end{longtable}
}

\begin{tcolorbox}[enhanced jigsaw, arc=.35mm, bottomrule=.15mm, bottomtitle=1mm, breakable, colback=white, colbacktitle=quarto-callout-warning-color!10!white, colframe=quarto-callout-warning-color-frame, coltitle=black, left=2mm, leftrule=.75mm, opacityback=0, opacitybacktitle=0.6, rightrule=.15mm, title=\textcolor{quarto-callout-warning-color}{\faExclamationTriangle}\hspace{0.5em}{伪重复（Pseudoreplication）}, titlerule=0mm, toprule=.15mm, toptitle=1mm]

把非独立的测量当作独立重复，是生态学论文被拒稿的最常见统计原因之一。Hurlbert
(1984) 的经典论文专门讨论这个问题。

\end{tcolorbox}

\section{常见保护类研究设计}\label{ux5e38ux89c1ux4fddux62a4ux7c7bux7814ux7a76ux8bbeux8ba1}

\subsection{完全随机设计（CRD）}\label{ux5b8cux5168ux968fux673aux8bbeux8ba1crd}

所有实验单位随机分配到各处理组。适用于同质性较高的实验材料。例如：温室中同一批次的幼苗随机分配三种药剂处理。

\subsection{随机区组设计（RCBD）}\label{ux968fux673aux533aux7ec4ux8bbeux8ba1rcbd}

当实验单位存在已知的系统性差异时（如坡位、土壤类型），先按差异因素分区组，再在区组内随机分配处理。

\subsection{野生动物调查设计}\label{ux91ceux751fux52a8ux7269ux8c03ux67e5ux8bbeux8ba1}

{\def\LTcaptype{none} % do not increment counter
\begin{longtable}[]{@{}lll@{}}
\toprule\noalign{}
方法 & 适用场景 & 关键参数 \\
\midrule\noalign{}
\endhead
\bottomrule\noalign{}
\endlastfoot
样线法 & 较大尺度的动物调查 & 样线长度、宽度、调查努力量 \\
样点法 & 鸟类调查 & 样点间距、观察时间、重复次数 \\
红外相机 & 野生动物监测 & 相机密度、工作时长、独立照片判定 \\
\end{longtable}
}

\section{常见错误}\label{ux5e38ux89c1ux9519ux8bef-1}

\begin{enumerate}
\def\labelenumi{\arabic{enumi}.}
\tightlist
\item
  只有1个对照样地却声称有对照
\item
  伪重复------把同一株的多片叶子当独立样本
\item
  把调查设计当实验设计做因果推断
\end{enumerate}

\section{练习}\label{ux7ec3ux4e60-2}

\begin{enumerate}
\def\labelenumi{\arabic{enumi}.}
\tightlist
\item
  判断以下哪个是真正的重复：

  \begin{itemize}
  \tightlist
  \item
    A. 同一株植物测5片叶子
  \item
    B. 5株植物各测1片叶子
  \item
    C. 同一个样方连续数5次昆虫\\
  \item
    D. 5个样方各调查一次
  \end{itemize}
\item
  设计一个''比较三种林型中鸟类丰富度''的调查方案，标注对照、重复和随机化方式
\end{enumerate}

\part{第二篇：看懂并整理保护数据}

\chapter{第4章
数据管理与可重复分析}\label{ux7b2c4ux7ae0-ux6570ux636eux7ba1ux7406ux4e0eux53efux91cdux590dux5206ux6790}

\chapter{第4章
数据管理与可重复分析}\label{ux7b2c4ux7ae0-ux6570ux636eux7ba1ux7406ux4e0eux53efux91cdux590dux5206ux6790-1}

\section{本章状态：seed}\label{ux672cux7ae0ux72b6ux6001seed}

本章内容尚在建设中。以下为本章规划大纲，具体内容将随课程反馈和新案例持续补充。

\section{真实保护问题}\label{ux771fux5b9eux4fddux62a4ux95eeux9898-3}

\section{学习目标}\label{ux5b66ux4e60ux76eeux6807-3}

\section{为什么这个问题需要统计}\label{ux4e3aux4ec0ux4e48ux8fd9ux4e2aux95eeux9898ux9700ux8981ux7edfux8ba1}

\section{核心概念}\label{ux6838ux5fc3ux6982ux5ff5}

\section{R 分析过程}\label{r-ux5206ux6790ux8fc7ux7a0b}

\begin{Shaded}
\begin{Highlighting}[]
\CommentTok{\# 本章 R 代码示例（待补充）}
\end{Highlighting}
\end{Shaded}

\section{结果解释}\label{ux7ed3ux679cux89e3ux91ca}

\section{常见错误}\label{ux5e38ux89c1ux9519ux8bef-2}

\section{AI 辅助任务}\label{ai-ux8f85ux52a9ux4efbux52a1-1}

\section{练习}\label{ux7ec3ux4e60-3}

\section{继续学习}\label{ux7ee7ux7eedux5b66ux4e60-1}

\chapter{第5章
描述统计与分布}\label{ux7b2c5ux7ae0-ux63cfux8ff0ux7edfux8ba1ux4e0eux5206ux5e03}

\chapter{第5章
描述统计与分布}\label{ux7b2c5ux7ae0-ux63cfux8ff0ux7edfux8ba1ux4e0eux5206ux5e03-1}

\begin{tcolorbox}[enhanced jigsaw, arc=.35mm, bottomrule=.15mm, bottomtitle=1mm, breakable, colback=white, colbacktitle=quarto-callout-note-color!10!white, colframe=quarto-callout-note-color-frame, coltitle=black, left=2mm, leftrule=.75mm, opacityback=0, opacitybacktitle=0.6, rightrule=.15mm, title=\textcolor{quarto-callout-note-color}{\faInfo}\hspace{0.5em}{本章状态：sprout}, titlerule=0mm, toprule=.15mm, toptitle=1mm]

已有正态分布核心内容（来自课程 wiki），需要补充描述统计完整内容。

\end{tcolorbox}

\section{认识你的数据}\label{ux8ba4ux8bc6ux4f60ux7684ux6570ux636e}

拿到 50
个样地的虫口密度数据，你首先想知道什么？------平均多少？变异多大？有没有异常值？

\section{正态分布：统计学最重要的分布}\label{ux6b63ux6001ux5206ux5e03ux7edfux8ba1ux5b66ux6700ux91cdux8981ux7684ux5206ux5e03}

12\textbar{} 13\textbar\# 正态分布（Normal Distribution） 14\textbar{}
15\textbar\#\# 定义 16\textbar{}
17\textbar 正态分布（又称高斯分布）是自然界中最常见、统计学中最重要的连续概率分布。由均值
μ 决定中心位置，由标准差 σ 决定离散程度。标准正态分布 μ=0, σ=1。
18\textbar{} 19\textbar\#\# 直观解释 20\textbar{}
21\textbar 如果你测量1000株同龄同种幼苗的高度，画成直方图------它会是一个中间高、两边低、左右对称的钟形。大多数值集中在均值附近，越远离均值的值越少见。
22\textbar{} 23\textbar\#\#\# 68-95-99.7 法则 24\textbar{} 25\textbar-
\textbf{±1σ}：约 68.27\% 的数据落在这个范围 26\textbar- \textbf{±2σ}：约
95.45\% 的数据落在这个范围------这是 α=0.05 的基础 27\textbar-
\textbf{±3σ}：约 99.73\% 的数据落在这个范围------超出即被视为异常值
28\textbar{} 29\textbar\#\# 在生物保护中的例子 30\textbar{} 31\textbar-
同一树种在相似生境中的胸径通常近似正态分布 32\textbar-
测量误差（如同一个人多次测量树高）通常服从正态分布 33\textbar-
大样本下的样本均值近似正态分布（中心极限定理） 34\textbar{}
35\textbar\#\# 为什么重要 36\textbar{}
37\textbar 许多统计方法（t检验、ANOVA、回归）都假设数据（或残差）服从正态分布。理解正态分布是理解这些方法的前提。
38\textbar{} 39\textbar\#\# 教学动画 40\textbar{} 41\textbar-
\href{../animation-scripts/normal-distribution-animation.md}{正态分布演示动画}
--- 10秒Canvas动画 42\textbar-
\href{../animation-scripts/normal-distribution-reference.md}{正态分布参考页}
--- 滚动式教学页面 43\textbar{} 44\textbar\#\# 相关概念 45\textbar{}
46\textbar- \href{descriptive-statistics.md}{描述性统计} 47\textbar-
\href{p-value.md}{P值} --- P值是在正态（或t）分布下计算的面积
48\textbar- \href{significance-level.md}{显著性水平} ---
α=0.05即正态分布双侧尾部各2.5\% 49\textbar{}

\section{R 分析过程}\label{r-ux5206ux6790ux8fc7ux7a0b-1}

\begin{Shaded}
\begin{Highlighting}[]
\CommentTok{\# 描述性统计}
\CommentTok{\# 待补充：均值、中位数、标准差、直方图}
\end{Highlighting}
\end{Shaded}

\section{练习}\label{ux7ec3ux4e60-4}

\chapter{第6章
样本、误差与不确定性}\label{ux7b2c6ux7ae0-ux6837ux672cux8befux5deeux4e0eux4e0dux786eux5b9aux6027}

\chapter{第6章
样本、误差与不确定性}\label{ux7b2c6ux7ae0-ux6837ux672cux8befux5deeux4e0eux4e0dux786eux5b9aux6027-1}

\section{本章状态：seed}\label{ux672cux7ae0ux72b6ux6001seed-1}

本章内容尚在建设中。以下为本章规划大纲，具体内容将随课程反馈和新案例持续补充。

\section{真实保护问题}\label{ux771fux5b9eux4fddux62a4ux95eeux9898-4}

\section{学习目标}\label{ux5b66ux4e60ux76eeux6807-4}

\section{为什么这个问题需要统计}\label{ux4e3aux4ec0ux4e48ux8fd9ux4e2aux95eeux9898ux9700ux8981ux7edfux8ba1-1}

\section{核心概念}\label{ux6838ux5fc3ux6982ux5ff5-1}

\section{R 分析过程}\label{r-ux5206ux6790ux8fc7ux7a0b-2}

\begin{Shaded}
\begin{Highlighting}[]
\CommentTok{\# 本章 R 代码示例（待补充）}
\end{Highlighting}
\end{Shaded}

\section{结果解释}\label{ux7ed3ux679cux89e3ux91ca-1}

\section{常见错误}\label{ux5e38ux89c1ux9519ux8bef-3}

\section{AI 辅助任务}\label{ai-ux8f85ux52a9ux4efbux52a1-2}

\section{练习}\label{ux7ec3ux4e60-5}

\section{继续学习}\label{ux7ee7ux7eedux5b66ux4e60-2}

\chapter{第7章
数据可视化}\label{ux7b2c7ux7ae0-ux6570ux636eux53efux89c6ux5316}

\chapter{第7章
数据可视化}\label{ux7b2c7ux7ae0-ux6570ux636eux53efux89c6ux5316-1}

\section{图表选择指南}\label{ux56feux8868ux9009ux62e9ux6307ux5357}

{\def\LTcaptype{none} % do not increment counter
\begin{longtable}[]{@{}lll@{}}
\toprule\noalign{}
想展示什么 & 推荐图表 & 例子 \\
\midrule\noalign{}
\endhead
\bottomrule\noalign{}
\endlastfoot
各组均值比较 & 柱状图+误差线 & 三种药剂的平均害虫数 \\
随时间变化趋势 & 折线图 & 连续5周的虫口密度 \\
数据分布形态 & 箱线图 & 不同林型的树高分布 \\
两个连续变量的关系 & 散点图 & 湿度与病害指数 \\
\end{longtable}
}

\section{箱线图怎么读}\label{ux7bb1ux7ebfux56feux600eux4e48ux8bfb}

箱线图展示了数据的五数概括：最小值、第一四分位数、中位数、第三四分位数、最大值。中间的粗线是中位数，箱子范围是中间50\%的数据。

\begin{tcolorbox}[enhanced jigsaw, arc=.35mm, bottomrule=.15mm, bottomtitle=1mm, breakable, colback=white, colbacktitle=quarto-callout-tip-color!10!white, colframe=quarto-callout-tip-color-frame, coltitle=black, left=2mm, leftrule=.75mm, opacityback=0, opacitybacktitle=0.6, rightrule=.15mm, title=\textcolor{quarto-callout-tip-color}{\faLightbulb}\hspace{0.5em}{Tip}, titlerule=0mm, toprule=.15mm, toptitle=1mm]

\textbf{为什么很多期刊不推荐柱状图？}
柱状图只展示均值（一个数字），丢失了数据分布的所有信息。相同的均值可能来自完全不同的数据分布。箱线图或散点图能让你看到数据的''形状''。

\end{tcolorbox}

\section{图表规范}\label{ux56feux8868ux89c4ux8303}

\begin{itemize}
\tightlist
\item
  \textbf{图题}放在图的下方，编号+描述
\item
  \textbf{坐标轴}必须标注变量名和单位
\item
  \textbf{误差线}必须说明是标准差(SD)还是标准误(SE)
\end{itemize}

\section{常见错误}\label{ux5e38ux89c1ux9519ux8bef-4}

\begin{enumerate}
\def\labelenumi{\arabic{enumi}.}
\tightlist
\item
  柱状图展示一切------分布形态应该用箱线图
\item
  误差线不标注含义------SD还是SE？
\item
  用3D图表------增加了视觉误导且难以精确读数
\end{enumerate}

\section{练习}\label{ux7ec3ux4e60-6}

同一组数据分别画柱状图、箱线图和散点图，比较哪种表达最合适。

\part{第三篇：回答常见保护问题}

\chapter{第8章 两组比较}\label{ux7b2c8ux7ae0-ux4e24ux7ec4ux6bd4ux8f83}

\chapter{第8章 两组比较}\label{ux7b2c8ux7ae0-ux4e24ux7ec4ux6bd4ux8f83-1}

\begin{tcolorbox}[enhanced jigsaw, arc=.35mm, bottomrule=.15mm, bottomtitle=1mm, breakable, colback=white, colbacktitle=quarto-callout-note-color!10!white, colframe=quarto-callout-note-color-frame, coltitle=black, left=2mm, leftrule=.75mm, opacityback=0, opacitybacktitle=0.6, rightrule=.15mm, title=\textcolor{quarto-callout-note-color}{\faInfo}\hspace{0.5em}{本章状态：sprout}, titlerule=0mm, toprule=.15mm, toptitle=1mm]

已有 P 值核心内容（来自课程 wiki），需要补充 t 检验完整流程和 R 代码。

\end{tcolorbox}

\section{真实保护问题}\label{ux771fux5b9eux4fddux62a4ux95eeux9898-5}

施药处理后，害虫数量真的变少了吗？还是只是巧合？

\section{P
值：统计推断的核心}\label{p-ux503cux7edfux8ba1ux63a8ux65adux7684ux6838ux5fc3}

11\textbar{} 12\textbar\# P 值（P-value） 13\textbar{} 14\textbar\#\#
定义 15\textbar{} 16\textbar P
值是\textbf{在零假设为真的前提下，观察到当前数据（或更极端数据）的概率}。
17\textbar{} 18\textbar- 它不是''零假设为真的概率'' 19\textbar-
它不是''处理无效的概率'' 20\textbar- 它不是''效应的大小'' 21\textbar{}
22\textbar\#\# 直观解释 23\textbar{} 24\textbar\textgreater{}
如果只让学生记住一句话：\textbf{P
值小，说明你看到的数据在''什么都没发生''的世界里不太可能出现------这让你怀疑''什么都没发生''可能不对。但
P 值不告诉你''到底发生了什么''、也不告诉你''效果有多大''。} 25\textbar{}
26\textbar 打个比方：你在森林里听到一声吼叫。如果假设''森林里没有老虎''，听到吼叫的
P 值就很小------你会怀疑这个假设。但 P
值不告诉你是一只老虎还是三只老虎，也不告诉你老虎离你多远。 27\textbar{}
28\textbar\#\# 在生物保护中的例子 29\textbar{}
30\textbar{}\textbf{病虫害防治试验}：你在 20 块样地中随机选 10
块施药，10 块做对照。施药后处理组平均虫口密度 3.2 头/株，对照组 5.8
头/株。t 检验得到 P = 0.03。 31\textbar{} 32\textbar- ✅ P = 0.03
意味着：如果药物真的完全无效（零假设为真），观察到这么大差异的概率只有
3\%。 33\textbar- ❌ P = 0.03 \textbf{不意味着}药物有效的概率是 97\%。
34\textbar- ❌ P = 0.03 \textbf{不意味着}药物能减少 2.6
头/株（这是效应量的工作）。 35\textbar{} 36\textbar\#\# 常见误解
37\textbar{} 38\textbar1. \textbf{``P \textless{} 0.05
就说明处理一定有效''} → 错。P
值只是反对零假设的证据强度，不能证明备择假设为真。且 P = 0.049 和 P =
0.051 之间没有本质差异。 39\textbar2. \textbf{``P 值越小，效应越大''} →
错。P 值受样本量影响极大。大样本下微小的效应也能得到很小的 P 值。
40\textbar3. \textbf{``P \textgreater{} 0.05 就说明没有效应''} →
错。可能是样本量太小（统计功效不足），而非真的没有效应。 41\textbar4.
\textbf{``P 值是零假设为真的概率''} → 错。这是最根本的误解。P
值是数据在零假设下出现的概率，不是零假设本身的概率。 42\textbar{}
43\textbar\#\# 相关概念 44\textbar{} 45\textbar-
\href{significance-level.md}{显著性水平} --- α = 0.05 从哪来？
46\textbar- \href{effect-size.md}{效应量} --- P 值不告诉你的东西
47\textbar- \href{statistical-power.md}{统计功效} --- 为什么 P
\textgreater{} 0.05 不等于无效 48\textbar-
\href{../concepts/type-i-error.md}{第一类错误} --- P \textless{} 0.05
仍然可能犯错 49\textbar-
\href{../concepts/confidence-interval.md}{置信区间} --- 比 P
值更丰富的信息 50\textbar{} 51\textbar\#\# 来源 52\textbar{} 53\textbar-
学生课堂反馈（2026-06-09）：P 值解释是高频误解来源 54\textbar{}

\section{R 分析过程}\label{r-ux5206ux6790ux8fc7ux7a0b-3}

\begin{Shaded}
\begin{Highlighting}[]
\CommentTok{\# 独立样本 t 检验示例}
\CommentTok{\# 待补充：处理组 vs 对照组害虫数量数据}
\end{Highlighting}
\end{Shaded}

\section{常见错误}\label{ux5e38ux89c1ux9519ux8bef-5}

\begin{enumerate}
\def\labelenumi{\arabic{enumi}.}
\tightlist
\item
  P \textless{} 0.05 就说明处理一定有效
\item
  P 值越小效应越大
\item
  P \textgreater{} 0.05 就说明没有效应
\end{enumerate}

\section{练习}\label{ux7ec3ux4e60-7}

\chapter{第9章
多组比较与方差分析}\label{ux7b2c9ux7ae0-ux591aux7ec4ux6bd4ux8f83ux4e0eux65b9ux5deeux5206ux6790}

\chapter{第9章
多组比较与方差分析}\label{ux7b2c9ux7ae0-ux591aux7ec4ux6bd4ux8f83ux4e0eux65b9ux5deeux5206ux6790-1}

\section{真实保护问题}\label{ux771fux5b9eux4fddux62a4ux95eeux9898-6}

你测试了三种药剂对病害指数的影响，每种药剂处理30株苗。药剂A平均病害指数2.1，药剂B为3.5，药剂C为4.8。三组之间有差异吗？哪种最好？

\section{为什么不能做三次t检验？}\label{ux4e3aux4ec0ux4e48ux4e0dux80fdux505aux4e09ux6b21tux68c0ux9a8c}

如果你做三次t检验（A vs B, B vs C, A vs
C），每次用α=0.05，三次都不犯第一类错误的概率是 0.95³ =
0.857。也就是说，你有\textbf{14.3\%的概率至少犯一次第一类错误}------远高于你以为的5\%。

方差分析（ANOVA）用一个F检验同时比较所有组，把第一类错误率控制在α水平。

\section{方差分析的基本思路}\label{ux65b9ux5deeux5206ux6790ux7684ux57faux672cux601dux8def}

ANOVA将总变异分解为两部分：

\begin{itemize}
\tightlist
\item
  \textbf{组间变异}：处理之间的差异
\item
  \textbf{组内变异}：同一处理组内部的个体差异
\end{itemize}

F = 组间变异 /
组内变异。如果F很大，说明组间差异远大于随机误差，处理效应可能存在。

\section{总体差异 vs
两两比较}\label{ux603bux4f53ux5deeux5f02-vs-ux4e24ux4e24ux6bd4ux8f83}

ANOVA的F检验只告诉你''至少有一组不同''，不告诉你是谁和谁不同。需要事后检验（Post-hoc
test，如Tukey HSD）来做两两比较。

\begin{tcolorbox}[enhanced jigsaw, arc=.35mm, bottomrule=.15mm, bottomtitle=1mm, breakable, colback=white, colbacktitle=quarto-callout-important-color!10!white, colframe=quarto-callout-important-color-frame, coltitle=black, left=2mm, leftrule=.75mm, opacityback=0, opacitybacktitle=0.6, rightrule=.15mm, title=\textcolor{quarto-callout-important-color}{\faExclamation}\hspace{0.5em}{Important}, titlerule=0mm, toprule=.15mm, toptitle=1mm]

\textbf{分析顺序}：先看F检验（总体差异是否显著）→
如果显著，再做两两比较。如果F不显著，不要做两两比较。

\end{tcolorbox}

\section{前提假设}\label{ux524dux63d0ux5047ux8bbe}

\begin{enumerate}
\def\labelenumi{\arabic{enumi}.}
\tightlist
\item
  各组数据独立
\item
  各组数据近似正态分布
\item
  各组方差齐性（Levene检验）
\end{enumerate}

\section{常见错误}\label{ux5e38ux89c1ux9519ux8bef-6}

\begin{enumerate}
\def\labelenumi{\arabic{enumi}.}
\tightlist
\item
  多组比较做多次t检验
\item
  F显著就不看效应量
\item
  忽略方差齐性假设
\end{enumerate}

\section{练习}\label{ux7ec3ux4e60-8}

比较三种防治方法对病害指数的影响。数据中有4组（含对照组），每组25个观测值。应该用什么方法？写出分析步骤。

\chapter{第10章
分类与比例数据}\label{ux7b2c10ux7ae0-ux5206ux7c7bux4e0eux6bd4ux4f8bux6570ux636e}

\chapter{第10章
分类与比例数据}\label{ux7b2c10ux7ae0-ux5206ux7c7bux4e0eux6bd4ux4f8bux6570ux636e-1}

\section{真实保护问题}\label{ux771fux5b9eux4fddux62a4ux95eeux9898-7}

两种杀菌剂处理后，药剂A的幼苗存活率85\%，药剂B为78\%。差异是真的还是碰巧？

这是比例数据的比较问题------不能用t检验。

\section{为什么比例数据不能当连续数据分析}\label{ux4e3aux4ec0ux4e48ux6bd4ux4f8bux6570ux636eux4e0dux80fdux5f53ux8fdeux7eedux6570ux636eux5206ux6790}

比例数据的三个特点： 1. 值域有界（0-1或0\%-100\%） 2.
方差依赖于均值（接近0\%或100\%时方差趋近于0） 3.
数据是''成功/失败''的计数，不是连续测量

\section{卡方检验}\label{ux5361ux65b9ux68c0ux9a8c}

卡方检验比较观测频数和期望频数。适用于分类变量之间的关联分析。

\textbf{适用条件}：期望频数小于5的格子不超过20\%。

\section{Fisher精确检验}\label{fisherux7cbeux786eux68c0ux9a8c}

当样本量小、期望频数不满足卡方条件时使用。不依赖大样本近似，直接计算精确概率。

\section{列联表}\label{ux5217ux8054ux8868}

{\def\LTcaptype{none} % do not increment counter
\begin{longtable}[]{@{}lllll@{}}
\toprule\noalign{}
处理 & 存活 & 死亡 & 合计 & 存活率 \\
\midrule\noalign{}
\endhead
\bottomrule\noalign{}
\endlastfoot
药剂A & 85 & 15 & 100 & 85\% \\
药剂B & 78 & 22 & 100 & 78\% \\
合计 & 163 & 37 & 200 & 81.5\% \\
\end{longtable}
}

\section{效应量：优势比和相对风险}\label{ux6548ux5e94ux91cfux4f18ux52bfux6bd4ux548cux76f8ux5bf9ux98ceux9669}

\begin{itemize}
\tightlist
\item
  \textbf{优势比(OR)}：药剂A的存活优势是B的多少倍
\item
  \textbf{相对风险(RR)}：药剂A相对于B的存活率比值
\end{itemize}

\section{常见错误}\label{ux5e38ux89c1ux9519ux8bef-7}

\begin{enumerate}
\def\labelenumi{\arabic{enumi}.}
\tightlist
\item
  把比例数据当连续数据，用t检验
\item
  忽略期望频数条件
\item
  卡方显著但不知道差异方向（要回到列联表看百分比）
\end{enumerate}

\section{练习}\label{ux7ec3ux4e60-9}

不同处理下植物发病率的比较------选择合适的方法并解释结果。

\chapter{第11章
相关与线性回归}\label{ux7b2c11ux7ae0-ux76f8ux5173ux4e0eux7ebfux6027ux56deux5f52}

\chapter{第11章
相关与线性回归}\label{ux7b2c11ux7ae0-ux76f8ux5173ux4e0eux7ebfux6027ux56deux5f52-1}

\section{相关分析}\label{ux76f8ux5173ux5206ux6790}

\subsection{相关系数}\label{ux76f8ux5173ux7cfbux6570}

Pearson相关系数 r 衡量两个连续变量之间线性关系的强度和方向。r ∈ {[}-1,
1{]}。

\begin{itemize}
\tightlist
\item
  r \textgreater{} 0：正相关（X增大，Y也增大）
\item
  r \textless{} 0：负相关（X增大，Y减小）
\item
  r ≈ 0：无线性相关
\end{itemize}

\subsection{相关 ≠ 因果}\label{ux76f8ux5173-ux56e0ux679c}

这是统计学最核心的教训之一。两个变量高度相关，不代表一个导致了另一个。

\begin{tcolorbox}[enhanced jigsaw, arc=.35mm, bottomrule=.15mm, bottomtitle=1mm, breakable, colback=white, colbacktitle=quarto-callout-warning-color!10!white, colframe=quarto-callout-warning-color-frame, coltitle=black, left=2mm, leftrule=.75mm, opacityback=0, opacitybacktitle=0.6, rightrule=.15mm, title=\textcolor{quarto-callout-warning-color}{\faExclamationTriangle}\hspace{0.5em}{Warning}, titlerule=0mm, toprule=.15mm, toptitle=1mm]

\textbf{经典反例}：冰激凌销量与溺水人数高度正相关。但吃冰激凌不会导致溺水------真正的原因是''夏季气温高''同时导致了冰激凌销量增加和更多人游泳。

\end{tcolorbox}

\textbf{在保护生物学中的陷阱}：发现游客数量与野生动物出现频次呈负相关→不能直接说''游客吓跑了动物''------也可能是游客多的季节恰好是动物繁殖期。

\section{简单线性回归}\label{ux7b80ux5355ux7ebfux6027ux56deux5f52}

回归分析用一个变量（自变量X）解释或预测另一个变量（因变量Y）的变化。

\textbf{回归系数β}：X每增加1个单位，Y平均变化β个单位。

\textbf{R²}：模型解释的变异比例。R²=0.7表示模型解释了Y的70\%变异。

\section{两种方法的区别}\label{ux4e24ux79cdux65b9ux6cd5ux7684ux533aux522b}

{\def\LTcaptype{none} % do not increment counter
\begin{longtable}[]{@{}lll@{}}
\toprule\noalign{}
& 相关分析 & 回归分析 \\
\midrule\noalign{}
\endhead
\bottomrule\noalign{}
\endlastfoot
目的 & 衡量关系强度 & 解释/预测 \\
方向 & 对称（X和Y地位相同） & 不对称（Y是被解释变量） \\
输出 & r & β, R², P值 \\
\end{longtable}
}

\section{常见错误}\label{ux5e38ux89c1ux9519ux8bef-8}

\begin{enumerate}
\def\labelenumi{\arabic{enumi}.}
\tightlist
\item
  相关显著就声称因果
\item
  回归显著但R²极低（\textless0.1）仍声称''有重要影响''
\item
  不检查残差就报告回归结果
\end{enumerate}

\section{练习}\label{ux7ec3ux4e60-10}

分析温度与虫口密度的关系。画出散点图，计算相关系数和回归线。

\chapter{第12章
从线性模型走向广义线性模型}\label{ux7b2c12ux7ae0-ux4eceux7ebfux6027ux6a21ux578bux8d70ux5411ux5e7fux4e49ux7ebfux6027ux6a21ux578b}

\chapter{第12章
从线性模型走向广义线性模型}\label{ux7b2c12ux7ae0-ux4eceux7ebfux6027ux6a21ux578bux8d70ux5411ux5e7fux4e49ux7ebfux6027ux6a21ux578b-1}

\section{本章状态：seed}\label{ux672cux7ae0ux72b6ux6001seed-2}

本章内容尚在建设中。以下为本章规划大纲，具体内容将随课程反馈和新案例持续补充。

\section{真实保护问题}\label{ux771fux5b9eux4fddux62a4ux95eeux9898-8}

\section{学习目标}\label{ux5b66ux4e60ux76eeux6807-5}

\section{为什么这个问题需要统计}\label{ux4e3aux4ec0ux4e48ux8fd9ux4e2aux95eeux9898ux9700ux8981ux7edfux8ba1-2}

\section{核心概念}\label{ux6838ux5fc3ux6982ux5ff5-2}

\section{R 分析过程}\label{r-ux5206ux6790ux8fc7ux7a0b-4}

\begin{Shaded}
\begin{Highlighting}[]
\CommentTok{\# 本章 R 代码示例（待补充）}
\end{Highlighting}
\end{Shaded}

\section{结果解释}\label{ux7ed3ux679cux89e3ux91ca-2}

\section{常见错误}\label{ux5e38ux89c1ux9519ux8bef-9}

\section{AI 辅助任务}\label{ai-ux8f85ux52a9ux4efbux52a1-3}

\section{练习}\label{ux7ec3ux4e60-11}

\section{继续学习}\label{ux7ee7ux7eedux5b66ux4e60-3}

\part{第四篇：保护专业应用}

\chapter{第13章
森林与植物保护研究}\label{ux7b2c13ux7ae0-ux68eeux6797ux4e0eux690dux7269ux4fddux62a4ux7814ux7a76}

\chapter{第13章
森林与植物保护研究}\label{ux7b2c13ux7ae0-ux68eeux6797ux4e0eux690dux7269ux4fddux62a4ux7814ux7a76-1}

\section{本章状态：seed}\label{ux672cux7ae0ux72b6ux6001seed-3}

本章内容尚在建设中。以下为本章规划大纲，具体内容将随课程反馈和新案例持续补充。

\section{真实保护问题}\label{ux771fux5b9eux4fddux62a4ux95eeux9898-9}

\section{学习目标}\label{ux5b66ux4e60ux76eeux6807-6}

\section{为什么这个问题需要统计}\label{ux4e3aux4ec0ux4e48ux8fd9ux4e2aux95eeux9898ux9700ux8981ux7edfux8ba1-3}

\section{核心概念}\label{ux6838ux5fc3ux6982ux5ff5-3}

\section{R 分析过程}\label{r-ux5206ux6790ux8fc7ux7a0b-5}

\begin{Shaded}
\begin{Highlighting}[]
\CommentTok{\# 本章 R 代码示例（待补充）}
\end{Highlighting}
\end{Shaded}

\section{结果解释}\label{ux7ed3ux679cux89e3ux91ca-3}

\section{常见错误}\label{ux5e38ux89c1ux9519ux8bef-10}

\section{AI 辅助任务}\label{ai-ux8f85ux52a9ux4efbux52a1-4}

\section{练习}\label{ux7ec3ux4e60-12}

\section{继续学习}\label{ux7ee7ux7eedux5b66ux4e60-4}

\chapter{第14章
野生动物调查与监测}\label{ux7b2c14ux7ae0-ux91ceux751fux52a8ux7269ux8c03ux67e5ux4e0eux76d1ux6d4b}

\chapter{第14章
野生动物调查与监测}\label{ux7b2c14ux7ae0-ux91ceux751fux52a8ux7269ux8c03ux67e5ux4e0eux76d1ux6d4b-1}

\section{本章状态：seed}\label{ux672cux7ae0ux72b6ux6001seed-4}

本章内容尚在建设中。以下为本章规划大纲，具体内容将随课程反馈和新案例持续补充。

\section{真实保护问题}\label{ux771fux5b9eux4fddux62a4ux95eeux9898-10}

\section{学习目标}\label{ux5b66ux4e60ux76eeux6807-7}

\section{为什么这个问题需要统计}\label{ux4e3aux4ec0ux4e48ux8fd9ux4e2aux95eeux9898ux9700ux8981ux7edfux8ba1-4}

\section{核心概念}\label{ux6838ux5fc3ux6982ux5ff5-4}

\section{R 分析过程}\label{r-ux5206ux6790ux8fc7ux7a0b-6}

\begin{Shaded}
\begin{Highlighting}[]
\CommentTok{\# 本章 R 代码示例（待补充）}
\end{Highlighting}
\end{Shaded}

\section{结果解释}\label{ux7ed3ux679cux89e3ux91ca-4}

\section{常见错误}\label{ux5e38ux89c1ux9519ux8bef-11}

\section{AI 辅助任务}\label{ai-ux8f85ux52a9ux4efbux52a1-5}

\section{练习}\label{ux7ec3ux4e60-13}

\section{继续学习}\label{ux7ee7ux7eedux5b66ux4e60-5}

\chapter{第15章
生物多样性与群落数据}\label{ux7b2c15ux7ae0-ux751fux7269ux591aux6837ux6027ux4e0eux7fa4ux843dux6570ux636e}

\chapter{第15章
生物多样性与群落数据}\label{ux7b2c15ux7ae0-ux751fux7269ux591aux6837ux6027ux4e0eux7fa4ux843dux6570ux636e-1}

\section{真实保护问题}\label{ux771fux5b9eux4fddux62a4ux95eeux9898-11}

你调查了两个样地的鸟类群落。样地A有12种鸟共240只，样地B也有12种鸟共240只。两个样地的多样性一样吗？

\textbf{不一定。}
如果样地A每种约20只（均匀分布），样地B有1种210只、其余11种共30只（极不均匀），那么样地A的多样性远高于样地B。

\section{核心指标}\label{ux6838ux5fc3ux6307ux6807}

\subsection{物种丰富度（Species Richness,
S）}\label{ux7269ux79cdux4e30ux5bccux5ea6species-richness-s}

样地中的物种数量。简单直观，但完全忽略了均匀度。

\subsection{Shannon指数（Shannon Index,
H'）}\label{shannonux6307ux6570shannon-index-h}

同时考虑物种数和均匀度。H'越高，多样性越高。

\subsection{均匀度（Evenness, J）}\label{ux5747ux5300ux5ea6evenness-j}

个体数量在各物种间的分布均匀程度。J=1表示完全均匀。

\section{关键注意事项}\label{ux5173ux952eux6ce8ux610fux4e8bux9879}

\subsection{调查努力量}\label{ux8c03ux67e5ux52aaux529bux91cf}

不同样地的调查努力量不同（如红外相机工作天数不同、样线长度不同），直接比较物种数会得出错误结论。需要使用稀疏曲线（Rarefaction
Curve）标准化。

\subsection{没有记录 ≠
不存在}\label{ux6ca1ux6709ux8bb0ux5f55-ux4e0dux5b58ux5728}

一次调查没有发现某个物种，不代表该物种不存在------可能只是未被检测到。这引出了检测概率（Detection
Probability）的概念。

\section{常见错误}\label{ux5e38ux89c1ux9519ux8bef-12}

\begin{enumerate}
\def\labelenumi{\arabic{enumi}.}
\tightlist
\item
  只看丰富度不看均匀度
\item
  调查努力量不同的样地直接比较
\item
  拿一次调查的物种数当绝对数字
\end{enumerate}

\section{练习}\label{ux7ec3ux4e60-14}

比较两个样地的鸟类多样性数据。计算丰富度、Shannon指数和均匀度。如果调查努力量不同，如何修正？

\chapter{第16章
措施效果与保护评估}\label{ux7b2c16ux7ae0-ux63aaux65bdux6548ux679cux4e0eux4fddux62a4ux8bc4ux4f30}

\chapter{第16章
措施效果与保护评估}\label{ux7b2c16ux7ae0-ux63aaux65bdux6548ux679cux4e0eux4fddux62a4ux8bc4ux4f30-1}

\section{本章状态：seed}\label{ux672cux7ae0ux72b6ux6001seed-5}

本章内容尚在建设中。以下为本章规划大纲，具体内容将随课程反馈和新案例持续补充。

\section{真实保护问题}\label{ux771fux5b9eux4fddux62a4ux95eeux9898-12}

\section{学习目标}\label{ux5b66ux4e60ux76eeux6807-8}

\section{为什么这个问题需要统计}\label{ux4e3aux4ec0ux4e48ux8fd9ux4e2aux95eeux9898ux9700ux8981ux7edfux8ba1-5}

\section{核心概念}\label{ux6838ux5fc3ux6982ux5ff5-5}

\section{R 分析过程}\label{r-ux5206ux6790ux8fc7ux7a0b-7}

\begin{Shaded}
\begin{Highlighting}[]
\CommentTok{\# 本章 R 代码示例（待补充）}
\end{Highlighting}
\end{Shaded}

\section{结果解释}\label{ux7ed3ux679cux89e3ux91ca-5}

\section{常见错误}\label{ux5e38ux89c1ux9519ux8bef-13}

\section{AI 辅助任务}\label{ai-ux8f85ux52a9ux4efbux52a1-6}

\section{练习}\label{ux7ec3ux4e60-15}

\section{继续学习}\label{ux7ee7ux7eedux5b66ux4e60-6}

\chapter{第17章
从统计结果到保护行动}\label{ux7b2c17ux7ae0-ux4eceux7edfux8ba1ux7ed3ux679cux5230ux4fddux62a4ux884cux52a8}

\chapter{第17章
从统计结果到保护行动}\label{ux7b2c17ux7ae0-ux4eceux7edfux8ba1ux7ed3ux679cux5230ux4fddux62a4ux884cux52a8-1}

\section{本章状态：seed}\label{ux672cux7ae0ux72b6ux6001seed-6}

本章内容尚在建设中。以下为本章规划大纲，具体内容将随课程反馈和新案例持续补充。

\section{真实保护问题}\label{ux771fux5b9eux4fddux62a4ux95eeux9898-13}

\section{学习目标}\label{ux5b66ux4e60ux76eeux6807-9}

\section{为什么这个问题需要统计}\label{ux4e3aux4ec0ux4e48ux8fd9ux4e2aux95eeux9898ux9700ux8981ux7edfux8ba1-6}

\section{核心概念}\label{ux6838ux5fc3ux6982ux5ff5-6}

\section{R 分析过程}\label{r-ux5206ux6790ux8fc7ux7a0b-8}

\begin{Shaded}
\begin{Highlighting}[]
\CommentTok{\# 本章 R 代码示例（待补充）}
\end{Highlighting}
\end{Shaded}

\section{结果解释}\label{ux7ed3ux679cux89e3ux91ca-6}

\section{常见错误}\label{ux5e38ux89c1ux9519ux8bef-14}

\section{AI 辅助任务}\label{ai-ux8f85ux52a9ux4efbux52a1-7}

\section{练习}\label{ux7ec3ux4e60-16}

\section{继续学习}\label{ux7ee7ux7eedux5b66ux4e60-7}

\part{第五篇：现代生态统计与继续学习}

\chapter{第18章
层级结构与混合模型}\label{ux7b2c18ux7ae0-ux5c42ux7ea7ux7ed3ux6784ux4e0eux6df7ux5408ux6a21ux578b}

\chapter{第18章
层级结构与混合模型}\label{ux7b2c18ux7ae0-ux5c42ux7ea7ux7ed3ux6784ux4e0eux6df7ux5408ux6a21ux578b-1}

\section{本章状态：seed}\label{ux672cux7ae0ux72b6ux6001seed-7}

本章内容尚在建设中。以下为本章规划大纲，具体内容将随课程反馈和新案例持续补充。

\section{真实保护问题}\label{ux771fux5b9eux4fddux62a4ux95eeux9898-14}

\section{学习目标}\label{ux5b66ux4e60ux76eeux6807-10}

\section{为什么这个问题需要统计}\label{ux4e3aux4ec0ux4e48ux8fd9ux4e2aux95eeux9898ux9700ux8981ux7edfux8ba1-7}

\section{核心概念}\label{ux6838ux5fc3ux6982ux5ff5-7}

\section{R 分析过程}\label{r-ux5206ux6790ux8fc7ux7a0b-9}

\begin{Shaded}
\begin{Highlighting}[]
\CommentTok{\# 本章 R 代码示例（待补充）}
\end{Highlighting}
\end{Shaded}

\section{结果解释}\label{ux7ed3ux679cux89e3ux91ca-7}

\section{常见错误}\label{ux5e38ux89c1ux9519ux8bef-15}

\section{AI 辅助任务}\label{ai-ux8f85ux52a9ux4efbux52a1-8}

\section{练习}\label{ux7ec3ux4e60-17}

\section{继续学习}\label{ux7ee7ux7eedux5b66ux4e60-8}

\chapter{第19章
生态过程与观测过程}\label{ux7b2c19ux7ae0-ux751fux6001ux8fc7ux7a0bux4e0eux89c2ux6d4bux8fc7ux7a0b}

\chapter{第19章
生态过程与观测过程}\label{ux7b2c19ux7ae0-ux751fux6001ux8fc7ux7a0bux4e0eux89c2ux6d4bux8fc7ux7a0b-1}

\section{本章状态：seed}\label{ux672cux7ae0ux72b6ux6001seed-8}

本章内容尚在建设中。以下为本章规划大纲，具体内容将随课程反馈和新案例持续补充。

\section{真实保护问题}\label{ux771fux5b9eux4fddux62a4ux95eeux9898-15}

\section{学习目标}\label{ux5b66ux4e60ux76eeux6807-11}

\section{为什么这个问题需要统计}\label{ux4e3aux4ec0ux4e48ux8fd9ux4e2aux95eeux9898ux9700ux8981ux7edfux8ba1-8}

\section{核心概念}\label{ux6838ux5fc3ux6982ux5ff5-8}

\section{R 分析过程}\label{r-ux5206ux6790ux8fc7ux7a0b-10}

\begin{Shaded}
\begin{Highlighting}[]
\CommentTok{\# 本章 R 代码示例（待补充）}
\end{Highlighting}
\end{Shaded}

\section{结果解释}\label{ux7ed3ux679cux89e3ux91ca-8}

\section{常见错误}\label{ux5e38ux89c1ux9519ux8bef-16}

\section{AI 辅助任务}\label{ai-ux8f85ux52a9ux4efbux52a1-9}

\section{练习}\label{ux7ec3ux4e60-18}

\section{继续学习}\label{ux7ee7ux7eedux5b66ux4e60-9}

\chapter{第20章
贝叶斯统计与模型诊断}\label{ux7b2c20ux7ae0-ux8d1dux53f6ux65afux7edfux8ba1ux4e0eux6a21ux578bux8bcaux65ad}

\chapter{第20章
贝叶斯统计与模型诊断}\label{ux7b2c20ux7ae0-ux8d1dux53f6ux65afux7edfux8ba1ux4e0eux6a21ux578bux8bcaux65ad-1}

\section{本章状态：seed}\label{ux672cux7ae0ux72b6ux6001seed-9}

本章内容尚在建设中。以下为本章规划大纲，具体内容将随课程反馈和新案例持续补充。

\section{真实保护问题}\label{ux771fux5b9eux4fddux62a4ux95eeux9898-16}

\section{学习目标}\label{ux5b66ux4e60ux76eeux6807-12}

\section{为什么这个问题需要统计}\label{ux4e3aux4ec0ux4e48ux8fd9ux4e2aux95eeux9898ux9700ux8981ux7edfux8ba1-9}

\section{核心概念}\label{ux6838ux5fc3ux6982ux5ff5-9}

\section{R 分析过程}\label{r-ux5206ux6790ux8fc7ux7a0b-11}

\begin{Shaded}
\begin{Highlighting}[]
\CommentTok{\# 本章 R 代码示例（待补充）}
\end{Highlighting}
\end{Shaded}

\section{结果解释}\label{ux7ed3ux679cux89e3ux91ca-9}

\section{常见错误}\label{ux5e38ux89c1ux9519ux8bef-17}

\section{AI 辅助任务}\label{ai-ux8f85ux52a9ux4efbux52a1-10}

\section{练习}\label{ux7ec3ux4e60-19}

\section{继续学习}\label{ux7ee7ux7eedux5b66ux4e60-10}

\chapter{第21章
时间、空间与随机过程}\label{ux7b2c21ux7ae0-ux65f6ux95f4ux7a7aux95f4ux4e0eux968fux673aux8fc7ux7a0b}

\chapter{第21章
时间、空间与随机过程}\label{ux7b2c21ux7ae0-ux65f6ux95f4ux7a7aux95f4ux4e0eux968fux673aux8fc7ux7a0b-1}

\section{本章状态：seed}\label{ux672cux7ae0ux72b6ux6001seed-10}

本章内容尚在建设中。以下为本章规划大纲，具体内容将随课程反馈和新案例持续补充。

\section{真实保护问题}\label{ux771fux5b9eux4fddux62a4ux95eeux9898-17}

\section{学习目标}\label{ux5b66ux4e60ux76eeux6807-13}

\section{为什么这个问题需要统计}\label{ux4e3aux4ec0ux4e48ux8fd9ux4e2aux95eeux9898ux9700ux8981ux7edfux8ba1-10}

\section{核心概念}\label{ux6838ux5fc3ux6982ux5ff5-10}

\section{R 分析过程}\label{r-ux5206ux6790ux8fc7ux7a0b-12}

\begin{Shaded}
\begin{Highlighting}[]
\CommentTok{\# 本章 R 代码示例（待补充）}
\end{Highlighting}
\end{Shaded}

\section{结果解释}\label{ux7ed3ux679cux89e3ux91ca-10}

\section{常见错误}\label{ux5e38ux89c1ux9519ux8bef-18}

\section{AI 辅助任务}\label{ai-ux8f85ux52a9ux4efbux52a1-11}

\section{练习}\label{ux7ec3ux4e60-20}

\section{继续学习}\label{ux7ee7ux7eedux5b66ux4e60-11}

\chapter{第22章 AI
辅助统计学习与分析}\label{ux7b2c22ux7ae0-ai-ux8f85ux52a9ux7edfux8ba1ux5b66ux4e60ux4e0eux5206ux6790}

\chapter{第22章 AI
辅助统计学习与分析}\label{ux7b2c22ux7ae0-ai-ux8f85ux52a9ux7edfux8ba1ux5b66ux4e60ux4e0eux5206ux6790-1}

\section{AI能做什么，不能做什么}\label{aiux80fdux505aux4ec0ux4e48ux4e0dux80fdux505aux4ec0ux4e48}

{\def\LTcaptype{none} % do not increment counter
\begin{longtable}[]{@{}ll@{}}
\toprule\noalign{}
能做什么 & 不能做什么 \\
\midrule\noalign{}
\endhead
\bottomrule\noalign{}
\endlastfoot
解释统计概念 & 判断方法是否适合你的具体数据 \\
推荐分析方法 & 替你承担分析错误的责任 \\
生成R/Python代码 & 保证代码可复现且正确 \\
解读输出结果 & 判断效应在管理上是否足够大 \\
\end{longtable}
}

\begin{quote}
\textbf{核心信念}：AI是工具，判断力是核心。
\end{quote}

\section{八要素提问模板}\label{ux516bux8981ux7d20ux63d0ux95eeux6a21ux677f}

向AI提问统计问题时，提供以下信息：

\begin{enumerate}
\def\labelenumi{\arabic{enumi}.}
\tightlist
\item
  我的研究问题是什么
\item
  我的研究对象是什么
\item
  我的变量有哪些
\item
  每个变量是什么类型
\item
  样本量是多少
\item
  分组情况是什么
\item
  我想比较差异、分析关系还是做预测
\item
  请给出适合的统计方法、理由和注意事项
\end{enumerate}

\begin{tcolorbox}[enhanced jigsaw, arc=.35mm, bottomrule=.15mm, bottomtitle=1mm, breakable, colback=white, colbacktitle=quarto-callout-tip-color!10!white, colframe=quarto-callout-tip-color-frame, coltitle=black, left=2mm, leftrule=.75mm, opacityback=0, opacitybacktitle=0.6, rightrule=.15mm, title=\textcolor{quarto-callout-tip-color}{\faLightbulb}\hspace{0.5em}{Tip}, titlerule=0mm, toprule=.15mm, toptitle=1mm]

\textbf{对比实验}：只给AI说''帮我分析这组数据''vs用八要素模板提问------得到的答案质量天差地别。

\end{tcolorbox}

\section{五点核查法}\label{ux4e94ux70b9ux6838ux67e5ux6cd5}

拿到AI回答后，必须检查：

{\def\LTcaptype{none} % do not increment counter
\begin{longtable}[]{@{}lll@{}}
\toprule\noalign{}
\# & 核查项 & 问题 \\
\midrule\noalign{}
\endhead
\bottomrule\noalign{}
\endlastfoot
1 & 数据类型 & AI正确识别了连续/计数/分类/比例吗？ \\
2 & 实验单位 & AI考虑了伪重复问题吗？ \\
3 & 方法选择 & AI推荐的方法适合你的设计和数据吗？ \\
4 & 代码可运行 & AI给的代码能在你的环境中运行吗？ \\
5 & 结论合理 & AI是否过度解释了结果？是否暗示了因果？ \\
\end{longtable}
}

\section{每次必须写一句人工核查}\label{ux6bcfux6b21ux5fc5ux987bux5199ux4e00ux53e5ux4ebaux5de5ux6838ux67e5}

格式：\textbf{``AI建议XXX，但经核查，XXX，因此我决定XXX。''}

\begin{quote}
示例：AI建议使用t检验，但本数据有三组处理，因此应改用方差分析。
\end{quote}

\section{AI常见错误类型}\label{aiux5e38ux89c1ux9519ux8befux7c7bux578b}

\begin{enumerate}
\def\labelenumi{\arabic{enumi}.}
\tightlist
\item
  \textbf{数据类型误判}：把计数型当连续型
\item
  \textbf{伪重复忽略}：不检查实验单位独立性
\item
  \textbf{过度因果解释}：从相关直接推断因果
\item
  \textbf{方法推荐不当}：对非正态数据仍推荐t检验
\item
  \textbf{忽视前提假设}：跳过正态性和方差齐性检验
\end{enumerate}

\section{练习}\label{ux7ec3ux4e60-21}

给你一组保护类数据，用八要素模板向AI提问，用五点核查法检查AI输出，写人工核查记录。

\part{附录}

\chapter{R 基础速查}\label{r-ux57faux7840ux901fux67e5}

\chapter{R 基础速查}\label{r-ux57faux7840ux901fux67e5-1}

\section{本附录状态：seed}\label{ux672cux9644ux5f55ux72b6ux6001seed}

内容尚在建设中。

\chapter{常用统计方法选择地图}\label{ux5e38ux7528ux7edfux8ba1ux65b9ux6cd5ux9009ux62e9ux5730ux56fe}

\chapter{常用统计方法选择地图}\label{ux5e38ux7528ux7edfux8ba1ux65b9ux6cd5ux9009ux62e9ux5730ux56fe-1}

\section{按研究目的选择}\label{ux6309ux7814ux7a76ux76eeux7684ux9009ux62e9}

{\def\LTcaptype{none} % do not increment counter
\begin{longtable}[]{@{}llll@{}}
\toprule\noalign{}
研究目的 & 因变量类型 & 自变量/分组 & 推荐方法 \\
\midrule\noalign{}
\endhead
\bottomrule\noalign{}
\endlastfoot
比较两组均值 & 连续型 & 两组 & 独立样本t检验 \\
比较多组均值 & 连续型 & 三组及以上 & 单因素方差分析(ANOVA) \\
比较比例 & 分类/比例 & 分组 & 卡方检验 \\
小样本比例 & 分类 & 分组 & Fisher精确检验 \\
分析两个连续变量的关系 & 连续型 & 连续型 & Pearson相关 / 线性回归 \\
计数型因变量 & 计数 & 分组/连续 & Poisson回归 \\
二分类因变量 & 0/1 & 分组/连续 & Logistic回归 \\
\end{longtable}
}

\section{按数据类型选择}\label{ux6309ux6570ux636eux7c7bux578bux9009ux62e9}

\begin{verbatim}
你的因变量是什么类型？
├── 连续型 → 自变量是什么类型？
│   ├── 两组分类 → t检验
│   ├── 多组分类 → ANOVA
│   └── 连续型 → 线性回归
├── 计数型 → Poisson回归 / 负二项回归
├── 分类/比例 → 自变量是什么类型？
│   ├── 分类 → 卡方检验 / Fisher精确检验
│   └── 连续型 → Logistic回归
└── 0/1 → Logistic回归
\end{verbatim}

\section{按研究设计选择}\label{ux6309ux7814ux7a76ux8bbeux8ba1ux9009ux62e9}

{\def\LTcaptype{none} % do not increment counter
\begin{longtable}[]{@{}ll@{}}
\toprule\noalign{}
设计类型 & 推荐分析方法 \\
\midrule\noalign{}
\endhead
\bottomrule\noalign{}
\endlastfoot
完全随机设计(CRD) & t检验 / 单因素ANOVA \\
随机区组设计(RCBD) & 双因素ANOVA / 线性混合模型 \\
前后对照 & 配对t检验 \\
重复测量 & 重复测量ANOVA / 线性混合模型 \\
样线/样点调查 & 根据变量类型选择，注意调查努力量标准化 \\
\end{longtable}
}

\section{常用非参数替代方法}\label{ux5e38ux7528ux975eux53c2ux6570ux66ffux4ee3ux65b9ux6cd5}

{\def\LTcaptype{none} % do not increment counter
\begin{longtable}[]{@{}lll@{}}
\toprule\noalign{}
参数方法 & 非参数替代 & 适用场景 \\
\midrule\noalign{}
\endhead
\bottomrule\noalign{}
\endlastfoot
独立t检验 & Mann-Whitney U检验 & 数据严重偏态 \\
配对t检验 & Wilcoxon符号秩检验 & 配对数据不满足正态 \\
单因素ANOVA & Kruskal-Wallis检验 & 多组非正态 \\
Pearson相关 & Spearman秩相关 & 非线性单调关系 \\
\end{longtable}
}

\chapter{统计术语中英文对照}\label{ux7edfux8ba1ux672fux8bedux4e2dux82f1ux6587ux5bf9ux7167}

\chapter{统计术语中英文对照}\label{ux7edfux8ba1ux672fux8bedux4e2dux82f1ux6587ux5bf9ux7167-1}

\section{本附录状态：seed}\label{ux672cux9644ux5f55ux72b6ux6001seed-1}

内容尚在建设中。

\chapter{可重复分析项目模板}\label{ux53efux91cdux590dux5206ux6790ux9879ux76eeux6a21ux677f}

\chapter{可重复分析项目模板}\label{ux53efux91cdux590dux5206ux6790ux9879ux76eeux6a21ux677f-1}

\section{本附录状态：seed}\label{ux672cux9644ux5f55ux72b6ux6001seed-2}

内容尚在建设中。

\chapter{数据伦理与敏感信息处理}\label{ux6570ux636eux4f26ux7406ux4e0eux654fux611fux4fe1ux606fux5904ux7406}

\chapter{数据伦理与敏感信息处理}\label{ux6570ux636eux4f26ux7406ux4e0eux654fux611fux4fe1ux606fux5904ux7406-1}

\section{本附录状态：seed}\label{ux672cux9644ux5f55ux72b6ux6001seed-3}

内容尚在建设中。

\chapter{AI
提问模板与核查清单}\label{ai-ux63d0ux95eeux6a21ux677fux4e0eux6838ux67e5ux6e05ux5355}

\chapter{AI
提问模板与核查清单}\label{ai-ux63d0ux95eeux6a21ux677fux4e0eux6838ux67e5ux6e05ux5355-1}

\section{八要素完整模板}\label{ux516bux8981ux7d20ux5b8cux6574ux6a21ux677f}

\begin{verbatim}
我的研究问题：[一句话描述]
研究对象：[物种/样地/处理]
变量列表：
  - [变量名]：[类型：连续/计数/分类/比例]，[角色：因变量/自变量]
  - ...
样本量：[每组/总体]
分组情况：[几组？每组多少？什么处理？]
分析目标：[比较差异 / 分析关系 / 做预测]
请给出：适合的统计方法、选择理由、前提假设和注意事项
\end{verbatim}

\section{五点核查清单}\label{ux4e94ux70b9ux6838ux67e5ux6e05ux5355}

{\def\LTcaptype{none} % do not increment counter
\begin{longtable}[]{@{}llll@{}}
\toprule\noalign{}
\# & 核查项 & ✓/✗ & 备注 \\
\midrule\noalign{}
\endhead
\bottomrule\noalign{}
\endlastfoot
1 & AI正确识别了数据类型吗？ & & \\
2 & AI考虑了实验单位和伪重复吗？ & & \\
3 & AI推荐的方法适合数据和研究设计吗？ & & \\
4 & AI给出的代码能运行吗？ & & \\
5 & AI是否过度解释了结果（暗示因果）？ & & \\
\end{longtable}
}

\section{人工核查记录模板}\label{ux4ebaux5de5ux6838ux67e5ux8bb0ux5f55ux6a21ux677f}

\begin{verbatim}
日期：
提问内容：[八要素]
AI回答摘要：
核查结果：
 - 数据类型：[✓/✗] [说明]
 - 实验单位：[✓/✗] [说明]
 - 方法选择：[✓/✗] [说明]
 - 代码：[✓/✗] [说明]
 - 结论：[✓/✗] [说明]
最终判断：[采用 / 修改后采用 / 不采用]
理由：
\end{verbatim}

\chapter{从基础统计到现代生态统计的学习路线}\label{ux4eceux57faux7840ux7edfux8ba1ux5230ux73b0ux4ee3ux751fux6001ux7edfux8ba1ux7684ux5b66ux4e60ux8defux7ebf}

\chapter{从基础统计到现代生态统计的学习路线}\label{ux4eceux57faux7840ux7edfux8ba1ux5230ux73b0ux4ee3ux751fux6001ux7edfux8ba1ux7684ux5b66ux4e60ux8defux7ebf-1}

\section{本附录状态：seed}\label{ux672cux9644ux5f55ux72b6ux6001seed-4}

内容尚在建设中。


\backmatter


\end{document}
